\section{Method}
\label{sec:method}

\subsection{Input format}
\label{sec:format}

\Cref{fig:intro} shows the same row in both formats, built from the training data described in \cref{sec:data}.
The options are shuffled each time a row is used, so the position of the correct option carries no information.

\paragraph{Decoder.}
The prompt follows LLM-as-Jev.
The user message holds the state, the question, the options numbered $[1]$ to $[K]$, and one instruction: select one option and answer only with its number.
The assistant turn is filled in up to ``Best answer: ['', and the model has to continue with the number of its choice and a closing bracket.
Qwen3's thinking mode is turned off.
Each option $k$ is therefore the token sequence of ``$k$]'', and the options form a trie over these sequences.

\paragraph{Encoder.}
The state and the question are the same.
The instruction asks to mark the correct option with ``yes'' and every other option with ``no'', and the options are listed as ``$[k]$ \texttt{[M]} description'', with the tokenizer's mask token in front of each description.
There is no chat template and nothing is generated.
One forward pass gives a vector of logits at every mask, and the score of option $k$ comes from the logits of the tokens ``yes'' and ``no'' at mask $k$, each taken with a leading space, as it appears after another word.

% Figure 1: the same row as a decoder (LLM-as-Jev) and as an encoder (ours). Colours: encoder / decoder / ink2.
\begin{figure}[t]
  \centering
  \begin{tikzpicture}[x=1in, y=1in,
      font=\sffamily\scriptsize,
      prompt/.style={draw=ink2, line width=0.4pt, rounded corners=2pt, fill=black!3, inner sep=4pt,
                     align=left, font=\ttfamily\scriptsize},
      box/.style={draw, rounded corners=2pt, inner sep=4pt, align=center, text width=2.85in},
      arr/.style={-{Stealth[length=4pt]}, line width=0.6pt, color=ink2},
      head/.style={font=\sffamily\bfseries\small, anchor=west}]
    % ---------------- decoder (left) ----------------
    \node[head, text=decoder] at (0, 0) {LLM-as-Jev: decoder};
    \node[prompt, text width=2.85in, anchor=north west] (dp) at (0, -0.15) {%
      State: What size does the screen need to be cut?\\
      Which solution better achieves the goal?\\
      Options:\\
      \relax[1] Cut to the size of the jar opening.\\
      \relax[2] Long enough to cover the whole jar.\\
      Select one option. Answer only with its number.\\
      \textbf{Best answer:} \relax[\,\textcolor{decoder}{\textbf{1\,]}}};
    \node[box, draw=decoder, below=0.22in of dp] (dt) {Loss on the \textbf{digits of the option number}\\
      (a trie), plus KL anchors to the base model};
    \node[box, draw=decoder, below=0.2in of dt] (dd) {Decision: the most probable option number};
    \draw[arr] (dp) -- (dt);  \draw[arr] (dt) -- (dd);
    % ---------------- encoder (right) ----------------
    \node[head, text=encoder] at (3.45, 0) {UltraBERT: encoder (this work)};
    \node[prompt, text width=2.2in, anchor=north west] (ep) at (3.45, -0.15) {%
      State: What size does the screen need to be cut?\\
      Which solution better achieves the goal?\\
      Mark the correct option with yes\\and every other option with no.\\
      Options:\\
      \relax[1] \colorbox{encoder!18}{[M]} Cut to the size of the jar opening.\\
      \relax[2] \colorbox{encoder!18}{[M]} Long enough to cover the whole jar.};
    \node[align=left, text=encoder] (s1) at ($(ep.south east)+(0.42in, 0.31in)$) {$s_1$};
    \node[align=left, text=encoder] (s2) at ($(ep.south east)+(0.42in, 0.14in)$) {$s_2$};
    \draw[arr, color=encoder] ($(ep.south east)+(0, 0.31in)$) -- (s1);
    \draw[arr, color=encoder] ($(ep.south east)+(0, 0.14in)$) -- (s2);
    \node[box, draw=encoder, text width=1.28in, anchor=north west] (sm) at ($(ep.south west)+(0,-0.22in)$) {%
      \textbf{Plain softmax}\\[1pt]
      $s_k=\text{logit}(\textit{yes})$\\[2pt]
      $p_k=\dfrac{e^{s_k}}{\sum_j e^{s_j}}$\\[3pt]
      loss $=-\log p_{\text{gold}}$};
    \node[box, draw=encoder, fill=encoder!8, text width=1.28in, anchor=north west] (sg) at ($(sm.north east)+(0.14in,0)$) {%
      \textbf{Sigmoid-norm}\\[1pt]
      $s_k=\text{logit}(\textit{yes})-\text{logit}(\textit{no})$\\[2pt]
      $p_k=\dfrac{\sigma(s_k)}{\sum_j\sigma(s_j)}$\\[3pt]
      loss $=-\log p_{\text{gold}}$};
    \node[box, draw=encoder, text width=2.7in, anchor=north west] (ed) at ($(sm.south west)+(0,-0.2in)$) {Decision: $\arg\max_k s_k$\\one forward pass, no generation};
    \draw[arr] ($(sm.north)+(0,0.22in)$) -- (sm.north);
    \draw[arr] ($(sg.north)+(0,0.22in)$) -- (sg.north);
    \draw[arr] (sm.south) -- (sm.south |- ed.north);
    \draw[arr] (sg.south) -- (sg.south |- ed.north);
  \end{tikzpicture}
  \caption{One row, two ways. The decoder reads the numbered options and writes the number of its answer; the
  loss is taken over the tokens of that number. The encoder gets a mask in front of each option and scores each
  option directly. Two of the three recipes are shown: the plain softmax over the options, which scores an
  option by the ``yes'' logit alone, and the sigmoid-normalised loss, which scores it by ``yes'' minus ``no''.
  Each option's score goes through a sigmoid, and the sum in the denominator keeps the options competing. The third recipe is the softmax over ``yes'' minus ``no''.}
  \label{fig:intro}
\end{figure}


\subsection{Loss functions}
\label{sec:loss}

\paragraph{Decoder: trie loss with KL anchors.}
At each node of the trie the model chooses among the tokens that are legal there.
The tree loss is the negative log-probability of the gold option, which is the product of the choices along its path, each normalised over the legal tokens only,
\begin{equation}
  \mathcal{L}_{\text{tree}} = -\sum_{t}\log\frac{\exp z_t(g_t)}{\sum_{v\in\mathcal{V}_t}\exp z_t(v)},
\end{equation}
where $z_t$ are the logits at step $t$, $g_t$ the gold token and $\mathcal{V}_t$ the legal tokens.
Three anchor terms keep the model close to its starting version $p_0$ \citep{li2026llmasjev}: $\mathcal{L}_{\text{mass}}$ keeps the total probability on the legal tokens, $\mathcal{L}_{\text{out}}$ keeps how the remaining probability is spread over the rest of the vocabulary, and $\mathcal{L}_{\text{pos}}$ keeps the next-token distribution at the start of the reply and after the closing bracket.
They are KL divergences to $p_0$, computed over its 64 most likely tokens and one bucket for the rest.
The loss is $\mathcal{L}=\mathcal{L}_{\text{tree}}+\lambda\,(\mathcal{L}_{\text{mass}}+\mathcal{L}_{\text{out}}+\mathcal{L}_{\text{pos}})$ with $\lambda=1$.

\paragraph{Encoder recipes.}
Let $\ell_k(v)$ be the logit of token $v$ at mask $k$, and let $g$ be the gold option.
A score $s_k$ says how strongly the model accepts option $k$; the decision is the option with the highest score.
Training raises the gold option's score relative to the other options, but no mask gets its own yes-or-no target.
Thus, a chosen option can still prefer ``no'': for example, scores of $-3$ and $-8$ select the first option.
We train the following four recipes:
\begin{itemize}
  \item \textbf{Yes-logit softmax.}
    The score is the ``yes'' logit alone, and the options compete through one softmax:
    \begin{equation}
      s_k=\ell_k(\textit{yes}),\qquad
      p_k=\frac{e^{s_k}}{\sum_{j=1}^{K}e^{s_j}},\qquad
      \mathcal{L}_{\text{softmax}}=-\log p_g.
    \end{equation}

  \item \textbf{Yes-minus-no softmax.}
    The score measures the preference for ``yes'' over ``no'', followed by the same softmax:
    \begin{equation}
      s_k=\ell_k(\textit{yes})-\ell_k(\textit{no}),\qquad
      p_k=\frac{e^{s_k}}{\sum_{j=1}^{K}e^{s_j}},\qquad
      \mathcal{L}_{\text{softmax}}=-\log p_g.
    \end{equation}

  \item \textbf{Yes-minus-no sigmoid normalisation.}
    Each score first goes through a sigmoid, and the resulting values are normalised by their sum:
    \begin{equation}
      s_k=\ell_k(\textit{yes})-\ell_k(\textit{no}),\qquad
      p_k=\frac{\sigma(s_k)}{\sum_{j=1}^{K}\sigma(s_j)},\qquad
      \mathcal{L}_{\text{sig-norm}}=-\log p_g.
    \end{equation}
    This has the form of the DeepSeek-V3 mixture-of-experts gate \citep{deepseekv3}, but the options play the role of experts and are normalised over all $K$ options, with no selection or balancing bias.
    Each $\sigma(s_k)$ lies between 0 and 1 and can be read as how strongly the model accepts option $k$ on its own, although only the ratios enter $p_k$.
    The options therefore still compete, and the decision remains the option with the highest score because the sigmoid is increasing.
    The gradient with respect to $s_k$ is
    \begin{equation}
      \left(1-\sigma(s_k)\right)
      \left(p_k-\mathbf{1}[k=g]\right).
    \end{equation}
    Once an option's sigmoid is close to 1, its score hardly changes.
    We do not pair this normalisation with the ``yes'' logit alone: vocabulary logits are around 200 (\cref{sec:precision}), so every sigmoid is close to 1, every option receives probability $1/K$, and the gradient is negligible.
    The ``yes''-minus-``no'' score is much smaller and avoids this saturation.
    Nothing in the loss makes $\sigma(s_k)$ a probability that holds on its own; \cref{sec:coupling} checks what the trained model does.

  \item \textbf{Yes-minus-no square-root softplus.}
    To test whether the sigmoid's upper bound matters, we use
    $h(s)=\sqrt{\log(1+e^s)}$:
    \begin{equation}
      s_k=\ell_k(\textit{yes})-\ell_k(\textit{no}),\qquad
      p_k=\frac{h(s_k)}{\sum_{j=1}^{K}h(s_j)},\qquad
      \mathcal{L}_{\text{sqrt-sp}}=-\log p_g.
    \end{equation}
    For a clearly rejected option ($s\ll0$), $h(s)\approx e^{s/2}$, so the options behave as in a softmax at twice the temperature.
    For a clearly accepted option ($s\gg0$), $h(s)\approx\sqrt{s}$ keeps growing, whereas the sigmoid stops at 1.
    Its score therefore continues to receive a small gradient where the sigmoid would receive almost none.
\end{itemize}
